% !TEX program = pdflatex
% Compile this handout from the repository root with make problem HW=hw02.
% Questions transcribed verbatim from Math Analysis I Assignment 2.pdf.
% Original wording, including "seqeunce" and "continuous fraction", is retained.
\documentclass[nohyper,nobib,nofonts]{tufte-handout}
% Shared packages and typography for the book and homework handouts.
\usepackage[T1]{fontenc}
% Give titletoc a separate part number before nameref wraps \@part.
\usepackage{etoolbox}
\makeatletter
\@ifclassloaded{tufte-book}{%
  \patchcmd{\@part}
    {\thepart\hspace{1em}}
    {\protect\numberline{\thepart}}
    {}
    {\PackageError{analysis-notes}{Could not format part contents labels}
      {Check the book class definition of \string\@part.}}
}{}
\makeatother
\usepackage{nameref}
\usepackage{amsmath,amssymb,amsthm}
\usepackage{graphicx}
\usepackage{booktabs}
% Packages used by the temporary Tufte specimen content.
\usepackage{fancyvrb,xspace,units,multicol,makeidx}
\fvset{fontsize=\small}
\usepackage{url}
\usepackage{lipsum} % Remove after all placeholder prose has been replaced.
\usepackage[backend=biber,natbib=true,style=numeric]{biblatex}
\addbibresource{bibliography.bib}

% Preserve the template's full citations in the margin.
\usepackage{xargs}
\renewcommandx{\cite}[3][1={0pt},2={}]{\sidenote[][#1]{\fullcite[#2]{#3}}}

\graphicspath{{figures/}}
\setkeys{Gin}{width=\linewidth,totalheight=\textheight,keepaspectratio}
\setcounter{secnumdepth}{1}

% Retain the borrowed title/contents design only for the book.
% Adapted from Kevin Godby's Tufte-LaTeX title/contents example:
% https://groups.google.com/forum/#!topic/tufte-latex/ujdzrktC1BQ
\makeatletter
\@ifclassloaded{tufte-book}{%
  \makeindex
  \setcounter{tocdepth}{0} % Parts and chapters only.
  % Title-only chapter headings; retain counters for sections and references.
  \titleformat{\chapter}[block]
    {\relax\ifthenelse{\NOT\boolean{@tufte@symmetric}}{\begin{fullwidth}}{}}
    {}
    {0pt}
    {\huge\rmfamily\itshape}
    [\ifthenelse{\NOT\boolean{@tufte@symmetric}}{\end{fullwidth}}{}]
  \renewcommand{\maketitlepage}{%
    \begingroup
    \setlength{\parindent}{0pt}
    {\fontsize{24}{24}\selectfont\textit{\@author}\par}
    \vspace{1.75in}
    \begin{fullwidth}
      \raggedright
      {\fontsize{36}{44}\selectfont\@title\par}
    \end{fullwidth}
    \vspace{0.5in}
    {\fontsize{14}{14}\selectfont\textsf{\smallcaps{\@date}}\par}
    \vfill
    {\fontsize{14}{14}\selectfont\textit{\@publisher}\par}
    \thispagestyle{empty}
    \endgroup
  }
  \titlecontents{part}
    [0pt]
    {\addvspace{0.75\baselineskip}\normalfont\normalsize}
    {\allcaps{Part~\thecontentslabel}\quad\allcaps}
    {\allcaps}
    {}
    [\vspace{0.5\baselineskip}]
  % Hide chapter numbers in the contents as well.
  \titlecontents{chapter}
    [4em]
    {\normalfont\normalsize}
    {\textit}
    {\textit}
    {\qquad\thecontentspage}
    [\vspace{0.5\baselineskip}]
}{}
\makeatother

% Shared notation and environments. Keep topic-specific macros local.
\newcommand{\N}{\mathbb{N}}
\newcommand{\Z}{\mathbb{Z}}
\newcommand{\Q}{\mathbb{Q}}
\newcommand{\R}{\mathbb{R}}
\newcommand{\C}{\mathbb{C}}

% Publication month and year used on the copyright page.
\newcommand{\monthyear}{%
  \ifcase\month\or January\or February\or March\or April\or May\or June\or
  July\or August\or September\or October\or November\or December\fi
  \space\number\year
}

% One shared counter makes results unambiguous within each section.
\theoremstyle{plain}
\newtheorem{theorem}{Theorem}[section]
\newtheorem{lemma}[theorem]{Lemma}
\newtheorem{proposition}[theorem]{Proposition}
\newtheorem{corollary}[theorem]{Corollary}
\theoremstyle{definition}
\newtheorem{definition}[theorem]{Definition}
\newtheorem{example}[theorem]{Example}
\newtheorem{exercise}[theorem]{Exercise}
\theoremstyle{remark}
\newtheorem{remark}[theorem]{Remark}

% BEGIN TUFTE PLACEHOLDER HELPERS
% Adapted from the supplied Tufte-LaTeX sample (Apache-2.0).
% Remove these documentation helpers once the specimens are replaced.
\newcommand{\hangp}[1]{\makebox[0pt][r]{(}#1\makebox[0pt][l]{)}}

%%
% Prints an asterisk that takes up no horizontal space.
% Useful in tabular environments.

%%
% Prints a trailing space in a smart way.


%%
% Some shortcuts for Tufte's book titles.  The lowercase commands will
% produce the initials of the book title in italics.  The all-caps commands
% will print out the full title of the book in italics.
\newcommand{\vdqi}{\textit{VDQI}\xspace}
\newcommand{\ei}{\textit{EI}\xspace}
\newcommand{\ve}{\textit{VE}\xspace}
\newcommand{\be}{\textit{BE}\xspace}
\newcommand{\VDQI}{\textit{The Visual Display of Quantitative Information}\xspace}
\newcommand{\EI}{\textit{Envisioning Information}\xspace}
\newcommand{\VE}{\textit{Visual Explanations}\xspace}
\newcommand{\BE}{\textit{Beautiful Evidence}\xspace}

\newcommand{\TL}{Tufte-\LaTeX\xspace}

% Typesets the font size, leading, and measure in the form of 10/12x26 pc.
\newcommand{\measure}[3]{#1/#2$\times$\unit[#3]{pc}}

% Macros for typesetting the documentation
\newcommand{\hlred}[1]{\textcolor{Maroon}{#1}}% prints in red
\newcommand{\hangleft}[1]{\makebox[0pt][r]{#1}}
\newcommand{\hairsp}{\hspace{1pt}}% hair space
\newcommand{\hquad}{\hskip0.5em\relax}% half quad space
\newcommand{\ie}{\textit{i.\hairsp{}e.}\xspace}
\newcommand{\eg}{\textit{e.\hairsp{}g.}\xspace}
\newcommand{\na}{\quad--}% used in tables for N/A cells
\providecommand{\XeLaTeX}{X\lower.5ex\hbox{\kern-0.15em\reflectbox{E}}\kern-0.1em\LaTeX}
\newcommand{\tXeLaTeX}{\XeLaTeX\index{XeLaTeX@\protect\XeLaTeX}}
% \index{\texttt{\textbackslash xyz}@\hangleft{\texttt{\textbackslash}}\texttt{xyz}}
\newcommand{\tuftebs}{\symbol{'134}}% a backslash in tt type in OT1/T1
\newcommand{\doccmdnoindex}[2][]{\texttt{\tuftebs#2}}% command name -- adds backslash automatically (and doesn't add cmd to the index)
\newcommand{\doccmddef}[2][]{%
  \hlred{\texttt{\tuftebs#2}}\label{template:cmd:#2}%
  \ifthenelse{\isempty{#1}}%
    {% add the command to the index
      \index{#2 command@\protect\hangleft{\texttt{\tuftebs}}\texttt{#2}}% command name
    }%
    {% add the command and package to the index
      \index{#2 command@\protect\hangleft{\texttt{\tuftebs}}\texttt{#2} (\texttt{#1} package)}% command name
      \index{#1 package@\texttt{#1} package}\index{packages!#1@\texttt{#1}}% package name
    }%
}% command name -- adds backslash automatically
\newcommand{\doccmd}[2][]{%
  \texttt{\tuftebs#2}%
  \ifthenelse{\isempty{#1}}%
    {% add the command to the index
      \index{#2 command@\protect\hangleft{\texttt{\tuftebs}}\texttt{#2}}% command name
    }%
    {% add the command and package to the index
      \index{#2 command@\protect\hangleft{\texttt{\tuftebs}}\texttt{#2} (\texttt{#1} package)}% command name
      \index{#1 package@\texttt{#1} package}\index{packages!#1@\texttt{#1}}% package name
    }%
}% command name -- adds backslash automatically
\newcommand{\docopt}[1]{\ensuremath{\langle}\textrm{\textit{#1}}\ensuremath{\rangle}}% optional command argument
\newcommand{\docarg}[1]{\textrm{\textit{#1}}}% (required) command argument
\newenvironment{docspec}{\begin{quotation}\small\ttfamily\parskip0pt\parindent0pt\ignorespaces}{\end{quotation}}% command specification environment
\newcommand{\docenv}[1]{\texttt{#1}\index{#1 environment@\texttt{#1} environment}\index{environments!#1@\texttt{#1}}}% environment name
\newcommand{\docenvdef}[1]{\hlred{\texttt{#1}}\label{template:env:#1}\index{#1 environment@\texttt{#1} environment}\index{environments!#1@\texttt{#1}}}% environment name
\newcommand{\docpkg}[1]{\texttt{#1}\index{#1 package@\texttt{#1} package}\index{packages!#1@\texttt{#1}}}% package name
\newcommand{\doccls}[1]{\texttt{#1}}% document class name
\newcommand{\docclsopt}[1]{\texttt{#1}\index{#1 class option@\texttt{#1} class option}\index{class options!#1@\texttt{#1}}}% document class option name
\newcommand{\docclsoptdef}[1]{\hlred{\texttt{#1}}\label{template:clsopt:#1}\index{#1 class option@\texttt{#1} class option}\index{class options!#1@\texttt{#1}}}% document class option name defined
\newcommand{\docmsg}[2]{\bigskip\begin{fullwidth}\noindent\ttfamily#1\end{fullwidth}\medskip\par\noindent#2}
\newcommand{\docfilehook}[2]{\texttt{#1}\index{file hooks!#2}\index{#1@\texttt{#1}}}
\newcommand{\doccounter}[1]{\texttt{#1}\index{#1 counter@\texttt{#1} counter}}
% END TUFTE PLACEHOLDER HELPERS

% Shared physical dimensions for interval and number-line diagrams.
\tikzset{
  interval diagram/.style={
    x=1cm,y=6mm,font=\footnotesize,
    color=black,line width=0.5pt,line cap=round,line join=round,
    >={Stealth[length=4pt,width=3pt]}
  },
  interval endpoint/.style={
    circle,draw=black,line width=0.5pt,
    minimum size=4pt,inner sep=0pt,outer sep=0pt
  },
  interval open/.style={interval endpoint,fill=white},
  interval closed/.style={interval endpoint,fill=black},
  interval label/.style={inner sep=1pt}
}

\allowdisplaybreaks[1]

% Photo order (duplicate/earlier shots are not repeated):
% Q1: 9069 (=9067), 9070 (includes 9068), 9071, 9072, 9073 top.
% Q2: 9073 lower portion, 9074.
% Q3: 9075, 9076, 9077, 9079.
% Q4: 9080, 9081, 9082, 9083, 9084.
% Q5: 9085, 9086, 9087, 9088.
% Q6: 9090, 9091, 9092, 9093, 9094, 9095, 9096, 9097.

% Match the assignment's question and subpart numbering.
\renewcommand{\thetheorem}{\arabic{theorem}}
\renewcommand{\labelenumi}{(\alph{enumi})}

\title{Homework 02}
\author{Lucas Barbosa, \texttt{lc6424}}
\date{}

\begin{document}
\maketitle

\begin{exercise}
  \label{ex:hw02:nested-bounded-intervals}
  Give an example of a sequence of nested nonempty bounded intervals
  $\{I_n\}$ in $\R$ such that $I_{n+1}\subset I_n$ for each
  $n=1,2,\ldots$ but
  \[
    \bigcap_{n=1}^{\infty} I_n=\varnothing.
  \]
  Explain what goes ``wrong''.
\end{exercise}

\begin{proof}[Solution]
  % !TEX root = ../hw02.tex
Let
\[
  I_n=\left(0,\frac1n\right),\qquad n\geq1.
\]
Each $I_n$ is a nonempty bounded interval because it contains $1/(2n)$
and is contained in $(0,1)$. Since $1/(n+1)<1/n$,
\[
  I_{n+1}=\left(0,\frac1{n+1}\right)
    \subset I_n.
\]
Thus the intervals are nested, as shown in
Figure~\ref{fig:hw02:open-intervals}.
\begin{marginfigure}
  \centering
  % !TEX root = ../../problems/hw02.tex
% (0,1/n), n = 1,2,3,..., drawn on the same horizontal scale.
\begin{tikzpicture}[interval diagram]
  \foreach \n/\endpoint in {1/1,2/{\tfrac12},3/{\tfrac13}} {
    \coordinate (left) at (0,{1-\n});
    \coordinate (right) at ({3.6/\n},{1-\n});
    \draw (left) -- (right);
    \node[interval open] at (left) {};
    \node[interval open] at (right) {};
    \node[interval label,above=3pt] at (right) {$\endpoint$};
  }
  \node[interval label,above=3pt] at (0,0) {$0$};
  \node at (1.8,-2.9) {$\vdots$};
\end{tikzpicture}

  \caption{The nested intervals $I_n=(0,1/n)$ exclude $0$.}
  \label{fig:hw02:open-intervals}
\end{marginfigure}

To prove that their intersection is empty, let $x\in\R$.
If $x\leq0$, then $x\notin I_1$. If $x>0$, the Archimedean
property gives a positive integer $n>1/x$.\sidenote{\textit{Archimedean property.}
  For every $M\in\R$, there is a positive integer $n>M$.}
Then $1/n<x$ and $x\notin I_n$.
Hence every real number is excluded from at least one interval,
and therefore
\[
  \bigcap_{n=1}^{\infty}I_n=\varnothing.
\]

The nested interval property guarantees a nonempty intersection
for nested nonempty bounded closed intervals in $\R$. Here the
missing hypothesis is \emph{closedness}. Every finite intersection
is nonempty, but the only point common to all the closures
$[0,1/n]$ is $0$, which is excluded from every $I_n$.

\end{proof}

\clearpage

\begin{exercise}
  \label{ex:hw02:nested-closed-intervals}
  Give an example of a sequence of nested nonempty closed intervals
  $\{I_n\}$ in $\R$ such that $I_{n+1}\subset I_n$ for each
  $n=1,2,\ldots$ but
  \[
    \bigcap_{n=1}^{\infty} I_n=\varnothing.
  \]
  Explain what goes ``wrong''.
\end{exercise}

\begin{proof}[Solution]
  % !TEX root = ../hw02.tex
Let
\[
  I_n=[n,\infty),\qquad n\geq1.
\]
Each $I_n$ is a nonempty closed interval in $\R$, since $n\in I_n$
and its complement $(-\infty,n)$ is open. Moreover,
\[
  I_{n+1}=[n+1,\infty)\subset[n,\infty)=I_n,
\]
so the intervals are nested; see
Figure~\ref{fig:hw02:unbounded-intervals}.
\begin{marginfigure}
  \centering
  % !TEX root = ../../problems/hw02.tex
% [n,infinity), n = 1,2,3,...; infinity is a direction, not an endpoint.
\begin{tikzpicture}[interval diagram]
  \foreach \n in {1,2,3} {
    \coordinate (left) at ({0.7*(\n-1)},{1-\n});
    \draw[->] (left) -- (3.6,{1-\n});
    \node[interval closed] at (left) {};
    \node[interval label,above=3pt] at (left) {$\n$};
  }
  \node[interval label,above=3pt] at (3.6,0) {$+\infty$};
  \node at (1.8,-2.9) {$\vdots$};
\end{tikzpicture}

  \caption{Each ray $[n,\infty)$ includes its finite left endpoint
    and extends indefinitely to the right.}
  \label{fig:hw02:unbounded-intervals}
\end{marginfigure}

For any $x\in\R$, the Archimedean property gives a positive
integer $n>x$.
Then $x\notin[n,\infty)=I_n$. Consequently, no real number belongs
to every $I_n$, and
\[
  \bigcap_{n=1}^{\infty}I_n=\varnothing.
\]

The missing hypothesis in the nested interval property is
\emph{boundedness}. All the intervals $I_n$ are unbounded, and
their left endpoints eventually exceed every fixed real number.

\end{proof}

\clearpage

\begin{exercise}
  \label{ex:hw02:nested-rational-intervals}
  Give an example of a sequence of nested nonempty bounded closed intervals
  $\{I_n\}$ in $\Q$ such that $I_{n+1}\subset I_n$ for each
  $n=1,2,\ldots$ but
  \[
    \bigcap_{n=1}^{\infty} I_n=\varnothing.
  \]
  Explain what goes ``wrong''.
\end{exercise}

\begin{proof}[Solution]
  % !TEX root = ../hw02.tex
For $n\geq1$, define
\[
  m_n=\lfloor10^n\sqrt2\rfloor,\qquad
  a_n=\frac{m_n}{10^n},\qquad b_n=\frac{m_n+1}{10^n},
\]
and let $I_n=[a_n,b_n]\cap\Q$. Since $m_n\in\Z$, both endpoints
are rational. Thus $I_n$ is nonempty (it contains $a_n$), bounded,
and closed relative to $\Q$, being the intersection of a closed
subset of $\R$ with $\Q$.

Because $\sqrt2$ is irrational, the definition of the floor function
gives $m_n<10^n\sqrt2<m_n+1$. Consequently,
\[
  a_n<\sqrt2<b_n,\qquad b_n-a_n=10^{-n}.
\]
Figure~\ref{fig:hw02:rational-nested-intervals} illustrates these
rational intervals around the excluded point $\sqrt2$.
\begin{marginfigure}
  \centering
  % !TEX root = ../../problems/hw02.tex
% Schematic [a_n,b_n] intersect Q; the irrational limit is never included.
\begin{tikzpicture}[interval diagram]
  \draw[densely dashed,black!45] (1.8,0.55) -- (1.8,-2.4);
  \node[interval label,above=3pt] at (1.8,0.55) {$\sqrt2$};
  \foreach \n/\left/\right in {1/0/3.6,2/0.6/3,3/1.2/2.4} {
    \draw (\left,{1-\n}) -- (\right,{1-\n});
    \node[interval closed] at (\left,{1-\n}) {};
    \node[interval closed] at (\right,{1-\n}) {};
    \node[interval open] at (1.8,{1-\n}) {};
    \node[interval label,above=3pt] at (\left,{1-\n}) {$a_\n$};
    \node[interval label,above=3pt] at (\right,{1-\n}) {$b_\n$};
  }
  \node at (1.8,-2.9) {$\vdots$};
\end{tikzpicture}

  \caption{Schematic intervals $[a_n,b_n]\cap\Q$.
    Hollow markers indicate that $\sqrt2\notin\Q$;
    all endpoints are rational.}
  \label{fig:hw02:rational-nested-intervals}
\end{marginfigure}

Multiplying by $10$ gives $10m_n<10^{n+1}\sqrt2<10m_n+10$.
Taking floors therefore yields
\[
  10m_n\leq m_{n+1}\leq10m_n+9.
\]
It follows that
\[
  a_{n+1}=\frac{m_{n+1}}{10^{n+1}}\geq a_n,
  \qquad
  b_{n+1}=\frac{m_{n+1}+1}{10^{n+1}}\leq b_n.
\]
Thus $I_{n+1}\subseteq I_n$. In fact, the inclusion is strict because
the lengths of the corresponding real intervals decrease by a
factor of $10$. By the density of $\Q$ in $\R$, any nonempty
open portion removed contains a rational number, so at least one
point of $I_n$ does not belong to $I_{n+1}$.

The inequalities
\[
  0<\sqrt2-a_n<10^{-n},\qquad
  0<b_n-\sqrt2<10^{-n}
\]
show that $a_n\to\sqrt2$ and $b_n\to\sqrt2$. If
$x\in\bigcap_{n=1}^{\infty}I_n$, then $x\in\Q$ and
$a_n\leq x\leq b_n$ for every $n$. Since both endpoint sequences
converge to $\sqrt2$, the squeeze theorem applied to the constant
sequence with value $x$ gives $x=\sqrt2$, contradicting $x\in\Q$.\sidenote{
  \textit{Squeeze theorem.} If $u_n\leq v_n\leq w_n$ and
  $u_n,w_n\to\ell$, then $v_n\to\ell$.}
Therefore
\[
  \bigcap_{n=1}^{\infty}I_n=\varnothing.
\]

What fails is the \emph{completeness of $\Q$}. The corresponding
closed intervals in $\R$ have intersection $\{\sqrt2\}$, but this
point is absent from $\Q$.

\end{proof}

\clearpage

\begin{exercise}
  \label{ex:hw02:decreasing-sequence}
  Consider the seqeunce
  \[
    x_n=\frac{1}{n^2}+\frac{2}{n^2}+\cdots+\frac{n}{n^2},
    \quad n=1,2,\ldots.
  \]
  \begin{enumerate}
    \item Show that it is monotone decreasing and bounded below.
    \item Show that it is Cauchy.
    \item Obtain the limit of $\{x_n\}$.
    \item Let $S=\{x_1,x_2,\ldots,x_n,\ldots\}$.
      Obtain both $\inf S$ and $\sup S$.
  \end{enumerate}
\end{exercise}

\begin{proof}[Solution]
  % !TEX root = ../hw02.tex
\leavevmode
\begin{enumerate}
  \item \textit{Monotonicity and a lower bound.}
    Using $1+2+\cdots+n=n(n+1)/2$, we obtain
    \[
      x_n=\frac1{n^2}\sum_{k=1}^n k
        =\frac{n(n+1)}{2n^2}
        =\frac12+\frac1{2n}.
    \]
    Hence, for every $n\geq1$,
    \[
      x_{n+1}-x_n
        =\frac1{2(n+1)}-\frac1{2n}
        =-\frac1{2n(n+1)}<0.
    \]
    Thus $(x_n)$ is strictly decreasing. Since $1/(2n)>0$,
    we also have $x_n>1/2$ for every $n\geq1$, so the sequence
    is bounded below by $1/2$.

  \item \textit{The Cauchy property.}
    Let $\varepsilon>0$. By the Archimedean property, choose
    $N\in\N$ such that $N>1/\varepsilon$. For all $m,n\geq N$,
    the triangle inequality gives
    \[
      \begin{aligned}
        |x_n-x_m|
          &=\left|\frac1{2n}-\frac1{2m}\right|\\
          &\leq\frac1{2n}+\frac1{2m}
          \leq\frac1N<\varepsilon.
      \end{aligned}
    \]
    Therefore $(x_n)$ is Cauchy.

  \item \textit{The limit.}
    Since $1/(2n)\to0$ as $n\to\infty$,
    \[
      \lim_{n\to\infty}x_n
        =\lim_{n\to\infty}\left(\frac12+\frac1{2n}\right)
        =\frac12.
    \]

  \item \textit{Infimum and supremum.}
    Part (a) shows that $1/2$ is a lower bound for $S$.
    If $c>1/2$, choose $n>1/(2(c-1/2))$. Then
    \[
      x_n=\frac12+\frac1{2n}<c,
    \]
    so $c$ is not a lower bound. Thus $\inf S=1/2$.
    Since $(x_n)$ is decreasing and $x_1=1\in S$, its maximum
    is $1$, and therefore $\sup S=1$.
\end{enumerate}

\end{proof}

\clearpage

\begin{exercise}
  \label{ex:hw02:factorial-sequence}
  Consider the sequence
  \[
    x_n=\sum_{k=0}^{n}\frac{1}{k!},\quad n=1,2,\ldots,
  \]
  in $\Q$ and define $S=\{x_1,x_2,\ldots,x_n,\ldots\}$ which is a subset of $\Q$.
  \begin{enumerate}
    \item Show that $S$ is a bounded set.
    \item Find $a=\inf S$ and $b=\sup S$.
    \item Is it true that $a$ and $b$ both lie in $\Q$?
      Explain why and why not.
  \end{enumerate}
\end{exercise}

\begin{proof}[Solution]
  % !TEX root = ../hw02.tex
\leavevmode
\begin{enumerate}
  \item \textit{Boundedness.}
    For $k\geq2$, each of the $k-1$ factors $2,3,\ldots,k$
    is at least $2$, so
    \[
      k!=1\cdot2\cdots k\geq2^{k-1}.
    \]
    Consequently, for $n\geq2$,
    \[
      \begin{aligned}
        2\leq x_n
          &=2+\sum_{k=2}^n\frac1{k!}\\
          &\leq2+\sum_{k=2}^n\frac1{2^{k-1}}\\
          &=2+\frac{\tfrac12(1-(\tfrac12)^{n-1})}{1-\tfrac12}
          =3-2^{1-n}<3.
      \end{aligned}
    \]
    Since $x_1=2$, it follows that $2\leq x_n<3$ for every
    $n\geq1$. Thus $S$ is bounded.

  \item \textit{Infimum and supremum.}
    We take the infimum and supremum in $\R$. Since $x_1=2$
    and all terms are at least $2$, we have $a=\inf S=2$.
    Moreover, $x_{n+1}-x_n=1/(n+1)!>0$, so $(x_n)$ is strictly
    increasing. Since part (a) gives an
    upper bound, the monotone convergence theorem implies that
    $x_n\to L$ for some $L\in\R$.\sidenote{\textit{Monotone convergence theorem.}
      Every increasing real sequence that is bounded above converges.}
    Each $x_n\leq L$. If $u<L$, convergence gives a term $x_n>u$,
    so $u$ is not an upper bound. Therefore $b=\sup S=L$.

    All derivatives of $\exp(x)$ are continuous and equal to
    $\exp(x)$, so their values at $0$ are $1$.
    Taylor's theorem at $0$, evaluated at $1$, therefore
    gives\sidenote{\textit{Taylor remainder.} For a function with
      $n+1$ continuous derivatives on $[0,1]$, its degree-$n$
      expansion about $0$, evaluated at $1$, has remainder
      $f^{(n+1)}(\xi)/(n+1)!$ for some $\xi\in(0,1)$.}
    \[
      e=\sum_{k=0}^n\frac1{k!}
        +\frac{\exp(\xi_n)}{(n+1)!},\qquad 0<\xi_n<1.
    \]
    The remainder satisfies
    \[
      0<e-x_n\leq\frac{e}{(n+1)!}\longrightarrow0.
    \]
    Hence $L=e$, and
    \[
      a=2,\qquad b=e=\sum_{k=0}^{\infty}\frac1{k!}.
    \]

  \item \textit{Rationality.}
    We have $a=2\in\Q$, but $b=e\notin\Q$. To prove the latter,
    suppose that $e=p/q$ with $p,q\in\N$. Choose an integer
    $N\geq\max\{q,2\}$. Then
    \[
      R=N!\left(e-\sum_{k=0}^N\frac1{k!}\right)
    \]
    is an integer because both $N!e=N!p/q$ and every $N!/k!$ with
    $0\leq k\leq N$ are integers. On the other hand,
    \[
      \begin{aligned}
        0<R
          &=\sum_{j=1}^{\infty}
            \frac1{(N+1)(N+2)\cdots(N+j)}\\
          &\leq\sum_{j=1}^{\infty}\frac1{(N+1)^j}
          =\frac1N\leq\frac12<1.
      \end{aligned}
    \]
    This contradicts $R\in\Z$. Thus $e$ is irrational.

    In particular, $S$ has no supremum in the ordered field $\Q$.
    Indeed, any rational upper bound $u$ must satisfy $u>e$.
    By the density of $\Q$ in $\R$, there is a rational $v$
    with $e<v<u$, and $v$ is a smaller upper bound for $S$.
\end{enumerate}

\end{proof}

\clearpage

\begin{exercise}
  \label{ex:hw02:newton-iteration}
  In class, we discussed the iteration process
  \[
    x_{n+1}=a+\frac{1}{x_n},\quad n=0,1,2,\ldots,
    \quad a\geq 1,\quad x_0>0,
  \]
  whose limit gives us the continuous fraction
  \[
    b=a+\frac{1}{a+\frac{1}{a+\frac{1}{a+\cdots}}}.
  \]
  That is, $b$ is a fixed point of the function
  \[
    f(x)=a+\frac{1}{x},\quad x>0,
  \]
  which is realized as a root of the function $F(x)=x-f(x)$ $(x>0)$.
  \begin{enumerate}
    \item Show that the iterative algorithm
      \[
        x_{n+1}=\frac{ax_n^2+2x_n}{1+x_n^2},
        \quad n=0,1,2,\ldots,\quad x_0>0,
      \]
      arises from Newton's method that may be used to approach $b$ of our problem.
    \item Assume that $\{x_n\}$ in (a) converges.
      Verify that the limit is indeed the continuous fraction $b$ defined above.
    \item Take $a=\frac{3}{2}$. Prove $0<x_n\leq 2$ for $n=1,2,\ldots$.
    \item Prove that $\{x_n\}_{n=1,2,\ldots}$ increases for any $x_0\ne 2$
      and is constant when $x_0=2$.
    \item Show that $\{x_n\}$ converges to $2$ with any initial state $x_0$
      (thereby establishing the ``global convergence'' of the algorithm).
  \end{enumerate}
\end{exercise}

\begin{proof}[Solution]
  % !TEX root = ../hw02.tex
\leavevmode
\begin{enumerate}
  \item \textit{Newton's method.}
    For $x>0$, let
    \[
      F(x)=x-a-\frac1x,\qquad F'(x)=1+\frac1{x^2}>0.
    \]
    Newton's method for $F(x)=0$ is therefore
    \[
      \begin{aligned}
        x_{n+1}
          &=x_n-\frac{F(x_n)}{F'(x_n)}\\
          &=x_n-\frac{x_n-a-1/x_n}{1+1/x_n^2}\\
          &=\frac{x_n(1+x_n^2)-(x_n^3-ax_n^2-x_n)}{1+x_n^2}\\
          &=\frac{ax_n^2+2x_n}{1+x_n^2},
      \end{aligned}
    \]
    which is the stated iteration. If $x_n>0$, then the numerator
    and denominator are positive. Since $x_0>0$, induction shows
    that every iterate is positive and the algorithm is well defined.

  \item \textit{Identification of the limit.}
    Suppose $x_n\to L$. Positivity implies $L\geq0$; we first
    rule out $L=0$. If $L=0$, there is an $N$ such that
    $0<x_n<1$ for all $n\geq N$. For these indices,
    \[
      \frac{x_{n+1}}{x_n}
        =\frac{ax_n+2}{1+x_n^2}>1.
    \]
    The tail is then increasing and satisfies $x_n\geq x_N>0$,
    contradicting $x_n\to0$. Hence $L>0$.

    Passing to the limit in the recurrence gives
    \[
      L=\frac{aL^2+2L}{1+L^2}.
    \]
    Thus $L(L^2-aL-1)=0$. Since $L>0$, we obtain
    \[
      L^2-aL-1=0,\qquad
      L=\frac{a+\sqrt{a^2+4}}2.
    \]
    This is the unique positive solution of $L=a+1/L$.
    The continued fraction $b$ is a positive fixed point of
    $f(x)=a+1/x$, so uniqueness implies $L=b$.

  \item \textit{Bounds when $a=3/2$.}
    For the remaining parts, set $a=3/2$ and write
    \[
      g(x)=\frac{\tfrac32x^2+2x}{1+x^2},\qquad x>0.
    \]
    We have $g(x)>0$ and
    \[
      2-g(x)
        =\frac{2(1+x^2)-\tfrac32x^2-2x}{1+x^2}
        =\frac{(x-2)^2}{2(1+x^2)}\geq0.
    \]
    The nonnegativity of the numerator is illustrated in
    Figure~\ref{fig:hw02:nonnegative-parabola}.
    \begin{marginfigure}
      \centering
      % !TEX root = ../../problems/hw02.tex
% IMG_9094: y = (x - 2)^2, with axes and intercepts made explicit.
\begin{tikzpicture}[x=0.6cm,y=0.38cm,>=Stealth,font=\footnotesize]
  \draw[<->,black!65] (-1.05,0) -- (5,0) node[right] {$x$};
  \draw[<->,black!65] (0,-0.85) -- (0,7) node[above] {$y$};
  \draw[<->,line width=0.65pt,domain=-0.55:4.55,samples=100,smooth]
    plot (\x,{(\x-2)^2});
  \fill (0,4) circle[radius=1.6pt];
  \fill (2,0) circle[radius=1.6pt];
  \node[below left,inner sep=2pt] at (0,0) {$0$};
  \node[left=3pt] at (0,4) {$4$};
  \node[below=3pt] at (2,0) {$2$};
\end{tikzpicture}

      \caption{The numerator $(x-2)^2$ is nonnegative,
        with equality only at $x=2$.}
      \label{fig:hw02:nonnegative-parabola}
    \end{marginfigure}
    Therefore $0<g(x)\leq2$ for every $x>0$, with $g(x)=2$
    if and only if $x=2$. Applying this to $x_n>0$ proves
    \[
      0<x_n\leq2\qquad\text{for every }n\geq1,
    \]
    regardless of whether $x_0\leq2$.

    % Keep the monotonicity proof and its margin plot on the same page.
    \newpage

  \item \textit{Monotonicity.}
    If $x_0=2$, then $g(2)=2$, so induction gives $x_n=2$
    for all $n\geq0$.

    Now suppose $x_0\ne2$. The equality condition in part (c)
    gives $0<x_1<2$. Repeatedly applying the same condition yields
    $0<x_n<2$ for every $n\geq1$. For these indices,
    \[
      \begin{aligned}
        x_{n+1}-x_n
          &=\frac{\tfrac32x_n^2+2x_n-x_n(1+x_n^2)}{1+x_n^2}\\
          &=\frac{x_n(-2x_n^2+3x_n+2)}{2(1+x_n^2)}\\
          &=\frac{x_n(2-x_n)(2x_n+1)}{2(1+x_n^2)}>0.
      \end{aligned}
    \]
    Every factor in the last expression is positive. Equivalently,
    the quadratic $-2x^2+3x+2=(2-x)(2x+1)$ is positive on
    $(-1/2,2)$, as shown in Figure~\ref{fig:hw02:concave-parabola}.
    \begin{marginfigure}
      \centering
      % !TEX root = ../../problems/hw02.tex
% IMG_9096: y = -2x^2 + 3x + 2, with axes and intercepts made explicit.
\begin{tikzpicture}[x=0.85cm,y=0.4cm,>=Stealth,font=\footnotesize]
  \draw[<->,black!65] (-1.3,0) -- (2.95,0) node[right] {$x$};
  \draw[<->,black!65] (0,-3.1) -- (0,4.2) node[above] {$y$};
  \draw[<->,line width=0.65pt,domain=-0.95:2.45,samples=100,smooth]
    plot (\x,{-2*(\x)*(\x)+3*(\x)+2});
  \foreach \xx/\yy in {-0.5/0,2/0,0/2} {
    \fill (\xx,\yy) circle[radius=1.6pt];
  }
  \node[below right,inner sep=2pt] at (0,0) {$0$};
  \node[above left=3pt] at (-0.5,0) {$-\tfrac12$};
  \node[above right=3pt] at (2,0) {$2$};
  \node[left=3pt] at (0,2) {$2$};
\end{tikzpicture}

      \caption{The quadratic $-2x^2+3x+2$ is positive between
        its roots $-1/2$ and $2$.}
      \label{fig:hw02:concave-parabola}
    \end{marginfigure}
    Thus $(x_n)_{n\geq1}$ is strictly increasing when $x_0\ne2$.
    The restriction $n\geq1$ is necessary because if $x_0>2$, the first
    iterate is below $2$ and hence below $x_0$.

  \item \textit{Global convergence.}
    If $x_0=2$, the sequence is constant and converges to $2$.
    If $x_0>0$ and $x_0\ne2$, parts (c) and (d) show that
    $(x_n)_{n\geq1}$ is increasing and bounded above by $2$.
    The monotone convergence theorem therefore gives a limit.
    By part (b), with $a=3/2$, this limit is
    \[
      \frac{\tfrac32+\sqrt{(\tfrac32)^2+4}}2
        =\frac{\tfrac32+\tfrac52}2=2.
    \]
    Adding the initial term does not affect convergence, so
    $x_n\to2$ for every allowed initial state $x_0>0$.
\end{enumerate}

\end{proof}

\end{document}
